\documentclass[10pt,a4paper]{amsart}
\usepackage{amsmath,amssymb,amsthm,graphicx}
% Retained preamble for the Stacks Project chapter snapshot in this directory.
% The chapter's own preamble.tex is a copy of the shared Stacks Project preamble; the block
% below is sliced MECHANICALLY, line for line, from it. Nothing was added, removed or
% redefined; the wrapper class, the xr-hyper/hyperref wiring and the externaldocument list
% are omitted because the assembler supplies the document class.
%
% lines taken from preamble.tex: 3, 5, 9, 11, 12, 156, 157, 158, ... (51 lines total)
%
% package coverage restored from the chapter's own preamble (2026-10-13): xy, xy-2cell

\usepackage[all]{xy}
\xyoption{2cell}
\UseAllTwocells

\newenvironment{reference}{\comment}{\endcomment}
\newenvironment{slogan}{\comment}{\endcomment}
\newenvironment{history}{\comment}{\endcomment}
\theoremstyle{plain}
\newtheorem{theorem}[subsection]{Theorem}
\newtheorem{proposition}[subsection]{Proposition}
\newtheorem{lemma}[subsection]{Lemma}
\theoremstyle{definition}
\newtheorem{definition}[subsection]{Definition}
\newtheorem{example}[subsection]{Example}
\newtheorem{exercise}[subsection]{Exercise}
\newtheorem{situation}[subsection]{Situation}
\theoremstyle{remark}
\newtheorem{remark}[subsection]{Remark}
\newtheorem{remarks}[subsection]{Remarks}
\numberwithin{equation}{subsection}
\def\lim{\mathop{\mathrm{lim}}\nolimits}
\def\colim{\mathop{\mathrm{colim}}\nolimits}
\def\Spec{\mathop{\mathrm{Spec}}}
\def\Hom{\mathop{\mathrm{Hom}}\nolimits}
\def\Ext{\mathop{\mathrm{Ext}}\nolimits}
\def\SheafHom{\mathop{\mathcal{H}\!\mathit{om}}\nolimits}
\def\SheafExt{\mathop{\mathcal{E}\!\mathit{xt}}\nolimits}
\def\Sch{\mathit{Sch}}
\def\Mor{\mathop{\mathrm{Mor}}\nolimits}
\def\Ob{\mathop{\mathrm{Ob}}\nolimits}
\def\Sh{\mathop{\mathit{Sh}}\nolimits}
\def\NL{\mathop{N\!L}\nolimits}
\def\CH{\mathop{\mathrm{CH}}\nolimits}
\def\proetale{{pro\text{-}\acute{e}tale}}
\def\etale{{\acute{e}tale}}
\def\QCoh{\mathit{QCoh}}
\def\Ker{\mathop{\mathrm{Ker}}}
\def\Im{\mathop{\mathrm{Im}}}
\def\Coker{\mathop{\mathrm{Coker}}}
\def\Coim{\mathop{\mathrm{Coim}}}
\DeclareMathSymbol{\boxtimes}{\mathbin}{AMSa}{"02}
\def\QCohstack{\mathcal{QC}\!\mathit{oh}}
\def\Cohstack{\mathcal{C}\!\mathit{oh}}
\def\Spacesstack{\mathcal{S}\!\mathit{paces}}
\def\Quotfunctor{\mathrm{Quot}}
\def\Hilbfunctor{\mathrm{Hilb}}
\def\Curvesstack{\mathcal{C}\!\mathit{urves}}
\def\Polarizedstack{\mathcal{P}\!\mathit{olarized}}
\def\Complexesstack{\mathcal{C}\!\mathit{omplexes}}
\def\Pic{\mathop{\mathrm{Pic}}\nolimits}
\def\Picardstack{\mathcal{P}\!\mathit{ic}}
\def\Picardfunctor{\mathrm{Pic}}
\def\Deformationcategory{\mathcal{D}\!\mathit{ef}}
\usepackage{cleveref}
\crefname{originallemma}{original lemma}{original lemmas}
\Crefname{originallemma}{Original lemma}{Original lemmas}
\makeatletter
\newlabel{algebraic-lemma-lift-morphism-presentations}{{04T1}{source}{}{sourcepointer-1}{}}
\newlabel{algebraic-lemma-lift-morphism-presentations@cref}{{[originallemma][0][]04T1}{[1][0][]source}{}{}{}}
\makeatother
\crefname{originallemma}{original lemma}{original lemmas}
\Crefname{originallemma}{Original lemma}{Original lemmas}
\makeatletter
\newlabel{lemma-points-cartesian}{{lemma-points-cartesian}{source}{}{sourcepointer-2}{}}
\newlabel{lemma-points-cartesian@cref}{{[originallemma][0][]lemma-points-cartesian}{[1][0][]source}{}{}{}}
\makeatother
\crefname{originallemma}{original lemma}{original lemmas}
\Crefname{originallemma}{Original lemma}{Original lemmas}
\makeatletter
\newlabel{lemma-property-spaces-too}{{lemma-property-spaces-too}{source}{}{sourcepointer-3}{}}
\newlabel{lemma-property-spaces-too@cref}{{[originallemma][0][]lemma-property-spaces-too}{[1][0][]source}{}{}{}}
\makeatother
\crefname{originallemma}{original lemma}{original lemmas}
\Crefname{originallemma}{Original lemma}{Original lemmas}
\makeatletter
\newlabel{spaces-morphisms-definition-surjective}{{03ME}{source}{}{sourcepointer-4}{}}
\newlabel{spaces-morphisms-definition-surjective@cref}{{[originallemma][0][]03ME}{[1][0][]source}{}{}{}}
\makeatother
\crefname{originallemma}{original lemma}{original lemmas}
\Crefname{originallemma}{Original lemma}{Original lemmas}
\makeatletter
\newlabel{spaces-morphisms-lemma-fppf-open}{{042S}{source}{}{sourcepointer-5}{}}
\newlabel{spaces-morphisms-lemma-fppf-open@cref}{{[originallemma][0][]042S}{[1][0][]source}{}{}{}}
\makeatother
\crefname{originallemma}{original lemma}{original lemmas}
\Crefname{originallemma}{Original lemma}{Original lemmas}
\makeatletter
\newlabel{spaces-morphisms-lemma-smooth-flat}{{04TA}{source}{}{sourcepointer-6}{}}
\newlabel{spaces-morphisms-lemma-smooth-flat@cref}{{[originallemma][0][]04TA}{[1][0][]source}{}{}{}}
\makeatother
\crefname{originallemma}{original lemma}{original lemmas}
\Crefname{originallemma}{Original lemma}{Original lemmas}
\makeatletter
\newlabel{spaces-morphisms-lemma-smooth-locally-finite-presentation}{{04AJ}{source}{}{sourcepointer-7}{}}
\newlabel{spaces-morphisms-lemma-smooth-locally-finite-presentation@cref}{{[originallemma][0][]04AJ}{[1][0][]source}{}{}{}}
\makeatother
\title{The sets of points of algebraic stacks carry a unique topology continuous for every morphism and open for flat lfp maps from spaces}
\author{The Stacks Project Authors}
\date{stacks-project:stacks-properties.tex:local-snapshot-20261012; GNU FDL 1.2 or later}
\begin{document}
\maketitle
\noindent\textit{Editorial note.} Statement span is lines 1154-1165 of sources/stacks-stacks-properties/pinned-original-source.tex; proof span is lines 1167-1252, the complete author proof (84 body lines).\par\medskip
\noindent\textit{Editorial note.} Navigation stubs: 7 cited-but-not-included labels are supplied as non-load-bearing pointers with their pinned tags where the snapshot tag index has them.\par\medskip
\noindent\textit{Editorial note.} extra\_packages is empty on purpose: the retained preamble already loads xy with options and a second declaration clashes in the assembler wrapper.\par\medskip
\noindent\textit{Editorial note.} Text is retained verbatim, including the author's own typos and placeholders.\par\medskip
\noindent\textit{Editorial note.} Queue index 0 of queue\_stacks-properties.json (pb=3301, stmt 1154-1165, proof 1167-1252).\par\medskip
\noindent\textit{Editorial note.} stacks-properties wave-1: fresh adjudication accept (2026-10-14); pool scanned while CONTROL.md is down (0 bytes) - §2 row deferred per Chief, claim reported to Chief 08:0x.\par\medskip
\noindent\textit{Editorial note.} Cross-chapter same-label warning (NOT an occupancy conflict, Chief ruling 2026-10-13 evening): label 'lemma-topology-points' also exists in stacks-spaces-properties — a different proposition in a different chapter file.\par\medskip
\noindent\textit{Allowed standard prerequisites (editorial).}\begin{itemize}
\item Explicit prerequisite boundary. The cited inputs are prerequisites; the argument in the retained span is not replaced by any cited result.
\end{itemize}\medskip
\ifdefined\hypertarget\hypertarget{sourcepointer-1}{}\fi
\noindent\textit{Source reference.} Original lemma 04T1: sources/stacks-stacks-properties/pinned-original-source.tex, source lines 1252--1252. Cited by the retained statement or proof but not part of the collected spans; supplied only as a navigation pointer. The argument does not depend on its content.\par\medskip
\ifdefined\hypertarget\hypertarget{sourcepointer-2}{}\fi
\noindent\textit{Source reference.} Original lemma lemma-points-cartesian: sources/stacks-stacks-properties/pinned-original-source.tex, source lines 1019--1036. Cited by the retained statement or proof but not part of the collected spans; supplied only as a navigation pointer. The argument does not depend on its content.\par\medskip
\ifdefined\hypertarget\hypertarget{sourcepointer-3}{}\fi
\noindent\textit{Source reference.} Original lemma lemma-property-spaces-too: sources/stacks-stacks-properties/pinned-original-source.tex, source lines 650--658. Cited by the retained statement or proof but not part of the collected spans; supplied only as a navigation pointer. The argument does not depend on its content.\par\medskip
\ifdefined\hypertarget\hypertarget{sourcepointer-4}{}\fi
\noindent\textit{Source reference.} Original lemma 03ME: sources/stacks-stacks-properties/pinned-original-source.tex, source lines 1252--1252. Cited by the retained statement or proof but not part of the collected spans; supplied only as a navigation pointer. The argument does not depend on its content.\par\medskip
\ifdefined\hypertarget\hypertarget{sourcepointer-5}{}\fi
\noindent\textit{Source reference.} Original lemma 042S: sources/stacks-stacks-properties/pinned-original-source.tex, source lines 1252--1252. Cited by the retained statement or proof but not part of the collected spans; supplied only as a navigation pointer. The argument does not depend on its content.\par\medskip
\ifdefined\hypertarget\hypertarget{sourcepointer-6}{}\fi
\noindent\textit{Source reference.} Original lemma 04TA: sources/stacks-stacks-properties/pinned-original-source.tex, source lines 1252--1252. Cited by the retained statement or proof but not part of the collected spans; supplied only as a navigation pointer. The argument does not depend on its content.\par\medskip
\ifdefined\hypertarget\hypertarget{sourcepointer-7}{}\fi
\noindent\textit{Source reference.} Original lemma 04AJ: sources/stacks-stacks-properties/pinned-original-source.tex, source lines 1252--1252. Cited by the retained statement or proof but not part of the collected spans; supplied only as a navigation pointer. The argument does not depend on its content.\par\medskip
\begin{lemma}
\label{lemma-topology-points}
There exists a unique topology on the sets of points
of algebraic stacks with the following properties:
\begin{enumerate}
\item for every morphism of algebraic stacks $\mathcal{X} \to \mathcal{Y}$
the map $|\mathcal{X}| \to |\mathcal{Y}|$ is continuous, and
\item for every morphism $U \to \mathcal{X}$ which is flat and locally
of finite presentation with $U$ an algebraic space
the map of topological spaces $|U| \to |\mathcal{X}|$ is continuous and open.
\end{enumerate}
\end{lemma}

\begin{proof}
Choose a morphism $p : U \to \mathcal{X}$ which is
surjective, flat, and locally of finite presentation
with $U$ an algebraic space. Such exist by the definition of an algebraic
stack, as a smooth morphism is flat and locally of finite presentation
(see
Morphisms of Spaces,
Lemmas \ref{spaces-morphisms-lemma-smooth-locally-finite-presentation} and
\ref{spaces-morphisms-lemma-smooth-flat}).
We define a topology on $|\mathcal{X}|$ by the rule:
$W \subset |\mathcal{X}|$ is open if and only if $|p|^{-1}(W)$ is open
in $|U|$. To show that this is independent of the choice of $p$, let
$p' : U' \to \mathcal{X}$ be another morphism which is surjective, flat,
locally of finite presentation from an algebraic space to
$\mathcal{X}$. Set $U'' = U \times_\mathcal{X} U'$
so that we have a $2$-commutative diagram
$$
\xymatrix{
U'' \ar[r] \ar[d] & U' \ar[d] \\
U \ar[r] & \mathcal{X}
}
$$
As $U \to \mathcal{X}$ and $U' \to \mathcal{X}$ are surjective, flat,
locally of finite presentation we see that $U'' \to U'$ and $U'' \to U$
are surjective, flat and locally of finite presentation, see
Lemma \ref{lemma-property-spaces-too}.
Hence the maps $|U''| \to |U'|$ and $|U''| \to |U|$ are continuous, open
and surjective, see
Morphisms of Spaces,
Definition \ref{spaces-morphisms-definition-surjective} and
Lemma \ref{spaces-morphisms-lemma-fppf-open}.
This clearly implies that our definition is independent of the choice
of $p : U \to \mathcal{X}$.

\medskip\noindent
Let $f : \mathcal{X} \to \mathcal{Y}$ be a morphism of algebraic stacks.
By
Algebraic Stacks, Lemma \ref{algebraic-lemma-lift-morphism-presentations}
we can find a $2$-commutative diagram
$$
\xymatrix{
U \ar[d]_x \ar[r]_a & V \ar[d]^y \\
\mathcal{X} \ar[r]^f & \mathcal{Y}
}
$$
with surjective smooth vertical arrows.
Consider the associated commutative diagram
$$
\xymatrix{
|U| \ar[d]_{|x|} \ar[r]_{|a|} & |V| \ar[d]^{|y|} \\
|\mathcal{X}| \ar[r]^{|f|} & |\mathcal{Y}|
}
$$
of sets. If $W \subset |\mathcal{Y}|$ is open, then by the definition
above this means exactly that $|y|^{-1}(W)$ is open in $|V|$. Since
$|a|$ is continuous we conclude that
$|a|^{-1}|y|^{-1}(W) = |x|^{-1}|f|^{-1}(W)$ is open in $|W|$ which means
by definition that $|f|^{-1}(W)$ is open in $|\mathcal{X}|$.
Thus $|f|$ is continuous.

\medskip\noindent
Finally, we have to show that if $U$ is an algebraic space, and
$U \to \mathcal{X}$ is flat and locally of finite presentation, then
$|U| \to |\mathcal{X}|$ is open. Let $V \to \mathcal{X}$ be surjective,
flat, and locally of finite presentation with $V$ an algebraic space.
Consider the commutative diagram
$$
\xymatrix{
|U \times_\mathcal{X} V| \ar[r]_e \ar[rd]_f &
|U| \times_{|\mathcal{X}|} |V| \ar[d]_c \ar[r]_d &
|V| \ar[d]^b \\
& |U| \ar[r]^a & |\mathcal{X}|
}
$$
Now the morphism $U \times_\mathcal{X} V \to U$ is surjective, i.e,
$f : |U \times_\mathcal{X} V| \to |U|$ is surjective.
The left top horizontal arrow is surjective, see
Lemma \ref{lemma-points-cartesian}.
The morphism $U \times_\mathcal{X} V \to V$ is flat and locally of finite
presentation, hence $d \circ e : |U \times_\mathcal{X} V| \to |V|$ is open,
see
Morphisms of Spaces, Lemma \ref{spaces-morphisms-lemma-fppf-open}.
Pick $W \subset |U|$ open. The properties above imply that
$b^{-1}(a(W)) = (d \circ e)(f^{-1}(W))$ is open, which by construction means
that $a(W)$ is open as desired.
\end{proof}

\end{document}
